\documentclass[11pt,a4paper]{article}
\usepackage[T1]{fontenc}
\usepackage[utf8]{inputenc}
\usepackage[czech]{babel}
\usepackage{lmodern,microtype}
\usepackage{amsmath,amssymb,mathtools}
\usepackage{geometry}
\geometry{left=19mm,right=19mm,top=21mm,bottom=21mm}
\usepackage[dvipsnames]{xcolor}
\usepackage{enumitem,tabularx,booktabs}
\usepackage[most]{tcolorbox}
\usepackage{fancyhdr,lastpage}
\usepackage[hidelinks]{hyperref}
\definecolor{deep}{HTML}{203955}
\definecolor{teal}{HTML}{176E76}
\definecolor{slate}{HTML}{506173}
\definecolor{pale}{HTML}{F3F7F9}
\setlength{\parindent}{0pt}
\setlength{\parskip}{5pt}
\setlist[enumerate]{leftmargin=6mm,itemsep=5pt,topsep=3pt}
\pagestyle{fancy}\fancyhf{}
\fancyhead[L]{\sffamily\footnotesize\color{slate} MAT · 4. ročník · kombinatorika}
\fancyhead[R]{\sffamily\footnotesize\color{slate} Hodina \LESSONNUM\, / 112}
\fancyfoot[L]{\sffamily\footnotesize\color{slate} Výukový podklad · samostatný dokument}
\fancyfoot[R]{\sffamily\footnotesize\color{slate} \thepage/\pageref{LastPage}}
\renewcommand{\headrulewidth}{0.35pt}
\renewcommand{\footrulewidth}{0pt}
\newcommand{\sect}[1]{\par\vspace{2mm}{\Large\sffamily\bfseries\color{deep} #1}\par\vspace{-0.5mm}\noindent\color{teal}\rule{\linewidth}{0.7pt}\color{black}\vspace{1mm}}
\tcbset{theobox/.style={enhanced,breakable,colback=pale,colframe=teal!45!white,boxrule=.6pt,arc=1mm,left=3mm,right=3mm,top=2mm,bottom=2mm}}
\newcommand{\LESSONNUM}{37}
\begin{document}

{\Huge\sffamily\bfseries\color{deep} Vlastnosti kombinačních čísel}\par
{\sffamily\large\color{teal} Výukový podklad · 37. vyučovací hodina}\par
\textbf{Hodinová dotace:} 45 minut\quad \textbf{Ročník:} 4. SPŠ\quad \textbf{Rozpis:} TP 2.2\par
\begin{tcolorbox}[theobox]\textbf{Výukový cíl.} Dokázat symetrii a Pascalův vztah, používat základní součtové identity a rozlišit jejich kombinatorické zdůvodnění.\end{tcolorbox}
\sect{Teorie a vysvětlení}
\textbf{Symetrie:} \(\displaystyle\binom nk=\binom n{n-k}\). Vybrat \(k\) osob znamená zároveň určit zbylých \(n-k\).\par
\textbf{Pascalův vztah:} Pro \(1\le k\le n-1\)
\[\binom nk=\binom{n-1}{k-1}+\binom{n-1}k.\]
Důkaz: při výběru \(k\)-členné skupiny rozdělíme případy na ty, které obsahují pevně zvolenou osobu P, a ty, které ji neobsahují. V prvním případě přidáme \(k-1\) z ostatních \(n-1\), ve druhém \(k\) z ostatních.\par
\textbf{Součet řádku:} \(\displaystyle\sum_{k=0}^n\binom nk=2^n\), neboť \(n\)-prvková množina má \(2^n\) podmnožin (u každého prvku volíme ano/ne).
\sect{Řešené příklady}
\textbf{Příklad 1.} Vypočti \(\binom{10}{8}\).\par
\textit{Řešení.} Symetrií \(\binom{10}{8}=\binom{10}{2}=45\).\par\smallskip
\textbf{Příklad 2.} Vypočti \(\binom94+\binom93\) jediným kombinačním číslem.\par
\textit{Řešení.} Pascalův vztah s \(n=10,k=4\): \(\binom94+\binom93=\binom{10}4=210\). Přímá kontrola: \(126+84=210\).\par\smallskip
\textbf{Příklad 3.} Vysvětli rovnost \(\binom{10}4=\binom94+\binom93\) na výběru čtyřčlenné komise.\par
\textit{Řešení.} V desetihlavé skupině vybereme konkrétní osobu P. Buď P není v komisi: \(\binom94\), nebo je: \(\binom93\). Případy se nepřekrývají a pokrývají všechny komise.\par\smallskip
\textbf{Příklad 4.} Kolik podmnožin má pětiprvková množina?\par
\textit{Řešení.} U každého z pěti prvků se nezávisle rozhodneme, zda patří do podmnožiny; \(2^5=32\). Současně je to \(\binom50+\cdots+\binom55\).\par\smallskip
\textbf{Příklad 5.} Zapiš \(\binom nk\) pomocí \(\binom{n-1}{k-1}\).\par
\textit{Řešení.} Z faktoriálů \(\binom nk=\frac{n}{k}\binom{n-1}{k-1}\), pro \(1\le k\le n\). Lze to ověřit po dosazení obou vzorců.\par\smallskip
\clearpage\sect{Procvičování pro studenty}
\begin{enumerate}
\item Zjednoduš \(\binom{12}{11}\).
\item Dopočti \(\binom73+\binom74\).
\item Vyjádři \(\binom83+\binom82\) jediným členem.
\item Urči součet \(\binom60+\binom61+\cdots+\binom66\).
\item Kolik výběrů ze sedmi prvků má lichý počet prvků?
\end{enumerate}
\sect{Řešení a zdůvodnění k procvičování}
\begin{enumerate}
\item Symetrie: \(\binom{12}1=12\).
\item Pascalův vztah: \(\binom73+\binom74=\binom84=70\). Také symetrií \(2\binom73=70\).
\item Pascalův vztah: \(\binom93=84\).
\item \(2^6=64\).
\item \(2^{7-1}=64\). Zdůvodnění: doplněním nebo odstraněním pevného prvku navzájem párujeme podmnožiny sudé a liché velikosti; celkem je \(2^7\) podmnožin.
\end{enumerate}
\sect{Metodické poznámky}
\textbf{Typické chyby:}
\begin{itemize}[leftmargin=5mm,itemsep=1pt]
\item Pascalův vztah má shodný horní index \(n-1\) na pravé straně.
\item \(\binom nk=\binom n{n-k}\) není obecně rovno \(\binom{k}n\).
\item Součet \(\sum\binom nk\) zahrnuje také prázdnou a úplnou množinu.
\end{itemize}
\textbf{Doporučený průběh 45 minut:} 0–12 symetrie a důkaz; 12–25 Pascalův vztah; 25–33 součet podmnožin; 33–42 úlohy; 42–45 shrnutí.\par
\textbf{Návaznost:} Následující hodina vizualizuje Pascalův vztah a propojí ho s počtem cest v mřížce.\par
\vspace{2mm}\noindent{\footnotesize\color{slate} \textbf{Východiska:} rozpis výuky MAT, 4. ročník, hodiny 29--41; tematický plán SPŠ; srovnání tematického rozsahu a náročnosti s realisticky.cz, oddíl Kombinatorika (lekce 090101--090119). Výklad, příklady a řešení jsou nově autorsky sestavené, nikoli převzatá sbírka.}\par
\end{document}